One limitation observed during the evaluation is that, under strong flooding attacks, the embedded device may become so constrained that it cannot provide timely feedback on event data. In particular, Prometheus queries may fail to deliver results because the detection process is not prioritized high enough in the system. To address this, future work will investigate methods to ensure that the detection algorithm always operates within a reserved portion of system resources that cannot be interrupted or consumed by other processes. This would guarantee reliable monitoring and responsiveness, even under conditions of severe resource exhaustion.\\

Furthermore, it will be necessary to test whether the detection system can also recognize attack types that were not explicitly part of the training dataset. Evaluating the model’s ability to generalize to unseen attacks will provide valuable insights into its robustness and practical applicability in dynamic real-world environments.\\

Another important step is the extension of the detection system beyond a single EVSE towards multiple EVSEs operating simultaneously. By integrating several charging stations into a distributed detection environment, robustness can be improved and anomalous behavior can be identified more effectively. Such a setup would not only increase resilience against individual failures but also strengthen the overall security of the charging infrastructure.\\

In addition, the LSTM model should be trained with a broader set of attack scenarios and extended feature sets. In particular, features capturing network traffic should be incorporated alongside hardware performance events. This would reduce the dependency on hardware-related signals alone and enable the system to detect anomalies reflected in communication patterns, thereby increasing detection accuracy and coverage.\\